\scanpage{109}\mspage{\handwritten{\unsure{9}}}

that \uline{this ambiguity however} is an \uline{irrelevant one.} To make this point a little clearer, I could say that a third mathematician, informed of the work of both, will readily think\unsure{.|-|\,}out a functor or rather a pair of functors, associating to any structure of Mr. X one of Mr. Y and conver\deleted{\unsure{s}}sely, in such a way that by composition of the two, we will associate to a X-structure \handwritten{(}$T$ say\handwritten{)}another $T'$, which \deleted{of course} will not be isomorphic to $T$ of course, but endowed with a canonical \deleted{\unsure{isom}} $\handwritten{\infty}$-equivalence (in the sense of Mr. X) $T\handwritten{\underset{\infty}{\simeq}}T'$, and the same on the Mr. Y side. Most probably, a fourth mathematician, faced with the same situation as the third, will get \deleted{a different}\inserted{his own} pair of functors to reconciliate Mr X and Mr Y, \deleted{but}\inserted{which} very probably \deleted{his pair of functors} wont be equivalent \deleted{to} \handwritten{(}I mean isomorphic\handwritten{)} to the previous one. Here however a fifth mathematician, informed about this new perplexity, will probably show that the two Y-structures $U$ and $U'$, associated by his two colleagues to \deleted{the} an X-structure $T$, while not isomorphic alas, admit however a canonical $\handwritten{\infty}$-equivalence between $U$ and $U'$\inserted{ (in the sense of the Y-theory)}. I could go on with a sixth mathematici\deleted{a}an, confronted with the same perplexity as the previous one, who winds up with another \inserted{\handwritten{$\infty$}-}equivalence between $U$ and $U'$ (without being informed of the work of the fifth), and a seventh reconciling them by discovering an $\handwritten{\infty}$-equivalence between these equivalences. The story of course is infinite, I better stop with seven mathematicians, a fair number nowadays to allow themselves getting involved with foundational matters \ldots\ There should be a mathematical statement though resuming in finite terms this infinite story, but in order to write it down I guess a \unsure{m}inimum amount of conceptual work, in the context of a given notion of $\handwritten{\infty}$-groupoids satisfying the desiderata (A) and (B), should be done, and I am by no means sure I will go through this, not in this letter anyhow.

\handwritten{10.} Now in the long last the explicit description I promised of $\underline{C}_{n+1}$ in terms of $\underline{C}_n$. As a matter of fact, I have a handful to propose ! O\deleted{\unsure{n}}ne choice, about the widest I \deleted{\unsure{c|w}}\unsure{w|c}ould think of, is: for every pair $(f,g)$ in $\underline{C}_n$ satisfying condition $(1)$ of (A), add one \deleted{arrow} new arrow $h$. To avoid set-theoretic difficulties though, we better \deleted{\unsure{reduce}} first modify the definition of $\underline{C}_{\unsure{o|0}}$ so that the set of its objects should be in the universe we are working in, preferably even it be countable. Or else and more reasonably, we will pic\handwritten{k} one $h$ for every isomorphism class of situations $(f,g)$ in $\underline{C}_n$. Another restriction to avoid too much redundancy -- this was the first definition \unsure{actu}ally that flipped to my mind the day before yesterday -- is to add a \uline{new} $h$ only when there is no ``old'' one, namely in $\underline{C}_n$, serving the same purpose. Then it came to my mind that there is a lot of redundancy still, thus there would be already in

% archive sheet number: 108
